\documentclass[10pt,a4paper]{article}
\usepackage[utf8]{inputenc}\usepackage[T1]{fontenc}\usepackage[english]{babel}
\usepackage{graphicx,float,booktabs,amsmath,siunitx,xcolor,tabularx}
\usepackage[margin=2.2cm]{geometry}\usepackage[hidelinks]{hyperref}
\hypersetup{pdfauthor={Thomas Schmiegelt},pdftitle={Two AI agents, one stellarator -- project report}}
\newcommand{\bild}[3]{\begin{figure}[H]\centering\includegraphics[width=#2\linewidth]{bilder/#1}\caption{#3}\end{figure}}
\newcommand{\zwei}[4]{\begin{figure}[H]\centering\includegraphics[width=0.49\linewidth]{bilder/#1}\hfill
  \includegraphics[width=0.49\linewidth]{bilder/#2}\caption{#3}\label{#4}\end{figure}}
\newcommand{\kasten}[2]{\par\smallskip\noindent\fcolorbox{gray}{gray!8}{\parbox{0.96\linewidth}{\textbf{#1} #2}}\par\smallskip}
\title{\textbf{Two AI Agents, One Stellarator}\\[2pt]\large Building an engineering simulation pipeline from an
  open dataset -- including the dead ends}
\author{Thomas Schmiegelt \textbar{} hobby project \texttt{stellarator-design}}
\date{25 October 2026}
\begin{document}\maketitle

\kasten{Hobby project, driven by curiosity.}{Not affiliated with any research institute or company, not peer
reviewed. All numbers are estimates from simplified models and may contain errors. The W7-X model is
\emph{inspired by} Wendelstein~7-X: coil paths come from public data \cite{simsopt}; casings, connectors, potting,
supports and wall thicknesses are my own simplified assumptions and do not describe the real device. Code,
computations and the draft text were produced with the help of AI agents; all figures are my own computations
and renderings.}

\begin{abstract}
A stellarator confines a fusion plasma with twisted 3D superconducting coils that carry large magnetic forces.
In four days, a human and two AI agents built an end-to-end pipeline from a public coil dataset to a
structural analysis of the coil casings: coil path $\to$ magnetic field $\to$ Lorentz force $\to$ casing
geometry $\to$ stress. This report is mostly about \emph{how} it was built: the division of labour between a local
agent (Hermes) and a coding agent (Claude Code), a first model that was wrong while looking right, a fully
parametric CAD model that broke the mesher, and a second, CAD-free code path that finally worked. For a
W7-X-like system of 70 coils the result is 99\,\% of casing stresses below 43\,MPa -- under strongly simplified
supports.
\end{abstract}

\tableofcontents

%======================================================================
\section{Setup}
\kasten{Stellarator in five sentences.}{A fusion plasma (over 100 million degrees) is held by a magnetic field
inside a ring-shaped vacuum vessel. To keep particles from drifting out, field lines must wind helically around
the ring. A tokamak gets this twist from a large current in the plasma itself; a stellarator gets it entirely
from external coils, which therefore have to be twisted in 3D. This makes steady-state operation easier but turns
each coil into a demanding 3D part with tight tolerances. The reference device is Wendelstein~7-X (IPP
Greifswald) with 50 non-planar and 20 planar superconducting coils \cite{bosch2013}.}

\subsection{Agents and responsibilities}
\begin{table}[H]\centering\small
\begin{tabularx}{\linewidth}{@{}lX@{}}\toprule
Human & Goal and direction, reviews results \emph{as images}, decides, supplies the key design ideas when stuck. \\
Hermes & Hermes Agent \cite{hermes}, run locally with open-weight models: literature research, project bookkeeping, writing reusable
  \emph{skills}, long-running jobs overnight. \\
Claude Code & Coding agent \cite{claudecode}: simulation pipeline, debugging, post-processing, figures. \\\bottomrule
\end{tabularx}\end{table}
The agents never talk directly. They coordinate through files in the repository: an append-only findings log
(\texttt{BEFUNDE.md}; every entry has observation, measurement, location, status), a handover file for decisions
and running jobs (\texttt{ERINNERUNGEN.md}), and skills -- short Markdown playbooks an agent loads before a task.
Long simulations are started detached (\texttt{setsid nohup}) so they survive a session ending or a context
compaction. The very first human intervention was a stop: Hermes started building tooling right away and was told
to do the research first.

\subsection{Tool chain}
\begin{table}[H]\centering\small
\begin{tabularx}{\linewidth}{@{}lX@{}}\toprule
Physics data & simsopt \cite{simsopt} (coil Fourier coefficients), DESC \cite{desc} (plasma shape optimisation) \\
Geometry & Python + gmsh/OpenCASCADE \cite{gmsh} (direct meshing), FreeCAD 1.1 headless \cite{freecad} (parametric CAD) \\
Solver & ElmerFEM 9.0 \cite{elmer}: magnetostatics and linear elasticity in one coupled run \\
Own code & Biot-Savart (numba), FFT field solver, field-line tracer, collision checks \\
Visualisation & PyVista/VTK \cite{pyvista}, ParaView \cite{paraview}, Blender 4.0 \cite{blender} via glTF \\
Hardware & one workstation, 32\,GB RAM, one GPU (reserved for the local agent) \\\bottomrule
\end{tabularx}\end{table}

%======================================================================
\section{Research phase}
Hermes collected the basics: W7-X has a major radius of 5.5\,m, five field periods and 2.5\,T on axis
\cite{bosch2013}. Each coil is wound from a cable-in-conduit conductor: 243 NbTi strands of 0.57\,mm, each with
144 filaments in copper, inside a co-extruded 16\,mm $\times$ 16\,mm aluminium jacket with 37\,\% void for
helium cooling \cite{rummel2005,maix2006}. From a set of 33 PDFs supplied by the user, four were relevant; the most
useful one sizes coil structures for EU-DEMO with a 660\,MPa stress limit and wedged inner legs
\cite{bachmann2024}. A rate-limited arXiv search added work on HTS strain sensitivity \cite{schmeling2025} and on
manufacturing tolerances (W7-X: 2\,mm) \cite{gil2025}.

\kasten{Lesson -- verify citations.}{The bibliography an agent produced during research contained entries
without authors and with wrong venues or years. For this report every reference was checked individually
against full text or publisher metadata (Crossref, arXiv API). Treat agent-written references as unverified
until checked.}

%======================================================================
\section{First results -- and why they were wrong}
\subsection{Validating the solver}
Before simulating coils, ElmerFEM was tested on a case with an exact answer: an infinite straight wire,
$B=\mu_0 I/2\pi r$. The first runs returned zero everywhere -- silently. Root causes: a mesh format the
converter misread, and wrong keywords in the solver input file, which Elmer ignores without an error. Starting
from Elmer's own regression test instead of hand-written input fixed it: 0.09\,\% error inside the wire, 2\,\%
outside (Fig.~\ref{fig:valid}).

\subsection{The surrogate model (report v0.1)}
Hermes believed the real coil paths were unavailable (the URLs it tried returned 404) and generated a
parametric surrogate: ellipses and bands, current calibrated to 2.5\,T. Compared with W7-X the ampere-turns per
coil matched to $0.90\times$. A report was written.
\zwei{v01_spulen_iso.png}{v01_feld_verteilung.png}{Report v0.1: the surrogate coil set and its field. Plausible
numbers -- but not a stellarator.}{fig:v01}

A topological check by the coding agent -- does each coil actually encircle the plasma axis? -- failed for 20 of
50 coils; the ``planar'' coils were horizontal bands above and below the plasma. The good match was a
coincidence. And the data had been available all along: simsopt ships the Fourier coefficients of the seven W7-X
base coils in its package data \cite{simsopt}. Two more bugs surfaced later and changed published numbers: a sign
error in the axis formula (calibration current 12.18\,kA instead of 13.49\,kA per turn), and Elmer choosing the
current direction itself, so 25 mirrored coils ran backwards (1.38\,T instead of 2.50\,T on axis). The latter is
now caught by an automated check.
\kasten{Lesson -- plausible is not correct.}{Agents are good at producing numbers that look right. Every geometry
gets an independent invariant check (winding number around the axis, clearances), and every field or force is
computed by a second, independent method before it goes into a report.}

%======================================================================
\section{Restart: mesh directly from the data}
With the real coil paths (7 base coils, 70 by symmetry) the expected physics appeared immediately: bean-shaped
flux surfaces and the characteristic island chain (Fig.~\ref{fig:feld}). A custom FFT voxel solver matched exact
Biot-Savart to 1.3\,\% but destroyed the flux surfaces during field-line tracing, because a Cartesian grid breaks
the device's symmetry -- so field lines are always traced with Biot-Savart.
\zwei{w7x_feldlinien_3d.png}{w7x_poincare_bewertung.png}{Field lines on flux surfaces (left), Poincaré section
and rotational transform (right).}{fig:feld}

The plan had been to model coils in CAD. Generating the mesh directly was faster and more robust: a rectangle
at every path point, two ruled half-surfaces per coil, meshed by gmsh. Sweeps and smooth lofts failed with
self-intersecting surfaces. A clearance check on centre lines replaced boolean intersection tests (1.5\,s
instead of more than 45\,min). Elmer with all 70 coils (3.3\,M tetrahedra, 680\,s, 10\,GB) matched Biot-Savart
on the axis to a median of 0.24\,\%.
\zwei{w7x_wickelpakete.png}{elmer_70_spulen_vergleich.png}{70 winding packs meshed straight from the path data
(left); FEM vs.\ Biot-Savart along a coil normal (right).}{fig:elmer}
\zwei{validierung_draht.png}{w7x_voxelsolver_validierung.png}{Validation: FEM on a straight wire (left), FFT
voxel solver vs.\ Biot-Savart (right).}{fig:valid}

A side study ran the pipeline backwards from a target of 1\,GW net electric power: plasma size and field
$\to$ own plasma shape (DESC) $\to$ coils (simsopt). With high-temperature superconductors the device ends up at
about 10\,m major radius with forces of roughly 1\,GN per coil -- about a thousand times W7-X. From here on, the
support structure is the real problem.
\zwei{uebersicht_1GW_A6_A45.png}{struktur_A6_HTS.png}{1\,GW sizing study (left) and a first support shell with
coupled field $\to$ force $\to$ stress analysis (right).}{fig:gw}

%======================================================================
\section{The parametric CAD branch}
A continuous support shell could not be constructed cleanly. The human proposed a technique from commercial CAD
(Creo's \emph{Variable Section Sweep}): per coil a closed guide curve as the casing's inner contour, sketches
normal to it, a loft through them; all walls grow outwards. The agent implemented it in headless FreeCAD as a
fully parametric model: spreadsheet parameters, sketch dimensions as expressions, symmetric coils as clones,
connectors between neighbours.
\zwei{freecad_gui_w7x.png}{gehaeuse_A45_HTS_verbunden.png}{Parametric FreeCAD model: W7-X-like casings with
connectors (left), compact 1\,GW design (right).}{fig:fc}

What went wrong:
\begin{enumerate}\itemsep1pt
\item \textbf{Estimated, not designed, walls.} The first coupled run worked but gave stresses up to eight times
  the allowable value.
\item \textbf{Wedges break the mesher.} Inner wedges (after \cite{bachmann2024}) produced a valid CAD solid, but
  gmsh failed with ``overlapping facets'': several boolean cuts on the same winding pack created coincident
  faces. Six attempts later (tolerances, fusing in FreeCAD, deliberate overlap; the last one hung gmsh for over an
  hour) the wedge was disabled.
\item \textbf{Collisions.} A loft bulges between sketch stations and corners of tilted neighbours overlap. An
  incremental assembly check (add coils one by one, count collisions, render a video) found about 100.
\item \textbf{Fix in CAD (human idea): subtract the neighbours.} Each casing minus its neighbours' envelopes plus
  a gap: 0 collisions in all 70 steps. Meshing the whole assembly stayed fragile.
\end{enumerate}
\zwei{struktur_A45_HTS_gehaeuse.png}{gehaeuse_A45_HTS_keil_oben.png}{First stress run of the CAD casings (left);
casings with inner wedges -- valid in CAD, not meshable (right).}{fig:keil}
\kasten{Lesson -- valid is not meshable.}{The more boolean operations meet on the same faces, the more sliver
faces appear, and the mesher fails far away from where the problem was created. Keep the analysis model separate
from the design model.}

%======================================================================
\section{The parallel branch: no CAD kernel}
The human's suggestion: \emph{``keep this morning's working approach for the analysis and run the CAD path in
parallel as an add-on.''} Each casing became two rectangular tubes on the exact coil path (outer contour and
winding-pack cavity) as separate surface meshes -- no boolean operations at all. Wall thicknesses follow from load
and available space; a disentangling step tests surface points of foreign casings (including intermediate rings
between stations) against the local cross-section and shrinks the cheapest wall: 4980 $\to$ 40 $\to$ 0 contact
points.

\begin{table}[H]\centering\small
\begin{tabular}{@{}ll@{}}\toprule
Model & steel casings ($E$ = 200\,GPa), homogenised winding packs ($E$ = 100\,GPa), linear elastic \\
Load & Lorentz force $\vec J\times\vec B$ from the magnetostatic solution in the same run \\
Supports & one clamped patch per casing (simplification; no connectors, no support ring) \\
Mesh / run time & 3.02\,M tetrahedra in 78\,s; solve 55\,min, about 14\,GB RAM \\
Force per coil & FEM 2.42\,MN vs.\ Biot-Savart 2.39\,MN (+1\,\%) \\
Casing von Mises & mean 11\,MPa, 99\,\% below 43\,MPa, max 110\,MPa; deformation below 1\,mm \\\bottomrule
\end{tabular}\end{table}
Because of the simplified supports and walls as thin as 3\,mm in places this is \emph{not} a strength proof --
but it is a reproducible end-to-end pipeline.
\zwei{aufbau_schritt35.png}{struktur_w7x_gehaeuse_netz.png}{Incremental collision check in the CAD branch, step
35/70 (left); von Mises stress from the CAD-free branch (right).}{fig:aufbau}
\zwei{w7x_feldlinien_spannung.png}{w7x_feldlinien_spannung_geschlossen.png}{Casing stress with magnetic field
lines on flux surfaces, cut open and closed.}{fig:fls}
\bild{freecad_und_elmer.png}{0.95}{Both branches side by side: parametric CAD (left), stress analysis from the
CAD-free pipeline (right).}

%======================================================================
\section{Visualisation: from mesh to Blender}
A 3D cut-away model shows each coil's layers: cable (NbTi strands in copper, helium), aluminium jacket, turn
insulation, ground insulation, potting (an assumption), steel casing; $12\times9=108$ conductors per coil. A
separate conductor model answers the question ``where does the current flow?'': almost entirely in the
superconducting NbTi; copper only takes over briefly during a quench; jacket and helium carry none.
\zwei{w7x_spule_querschnitt.png}{w7x_torus3d_detail.png}{Coil layers (schematic, left) and in the 3D model
(right).}{fig:quer}
\zwei{w7x_leiter3d.png}{w7x_leiter3d_blender.png}{Conductor: opened jacket, twisted strands, one strand
magnified with NbTi filaments (count symbolic). PyVista (left), Blender Cycles (right).}{fig:leiter}
\bild{w7x_torus3d_blender.png}{0.95}{The full coil system with field lines, rendered in Blender.}

Getting PyVista's glTF export into Blender needed three fixes, each found by measuring instead of guessing:
colours arrived as sRGB values in a linear slot (washed out; fixed by linearising), the scene was rotated by
$90^\circ$ about the vertical axis (fixed at the root node), and the distribution's Blender build lacks the
denoiser (render fails; denoising disabled). Rendering runs on the CPU because the GPU is reserved for the local
agent.

%======================================================================
\section{Takeaways for working with coding agents}
\begin{itemize}\itemsep1pt
\item \textbf{Verification is the job.} The most expensive error looked plausible. Independent checks
  (invariants, second methods, regression tests from upstream projects) caught every major bug.
\item \textbf{Silent failure is the default in scientific tools.} Wrong keywords ignored, current directions
  chosen by the solver, meshers hanging for an hour. Every pipeline stage now asserts on its output.
\item \textbf{Prefer the boring path.} Direct mesh generation from data beat a CAD kernel with booleans. Keeping a
  working path alive while exploring the elegant one saved the project.
\item \textbf{Humans supplied the decisive ideas} (research first, variable section sweep, neighbour subtraction,
  the parallel branch); agents supplied speed in implementing, testing and documenting them.
\item \textbf{Files are the shared memory.} A findings log, a handover file and skills let two agents work on one
  project across sessions and context compactions.
\item \textbf{References need checking} like any other agent output.
\end{itemize}

\subsection*{Open items}
Potting layer between winding pack and casing; connectors and support ring instead of clamped patches; minimum
wall thickness; wedge and casing in one body; wall sizing for the 1\,GW design; HTS strain check; tolerance study
($\pm$2\,mm).

\subsection*{Code and data}
\url{https://github.com/ThomasSchmiegelt/stellarator-design} (MIT licence): pipeline code, findings log, figures.

\begin{thebibliography}{99}\small
\bibitem{bosch2013} H.-S. Bosch, R.\,C. Wolf, T. Andreeva et al.: Technical challenges in the construction of
  the steady-state stellarator Wendelstein 7-X. \emph{Nucl.\ Fusion} 53 (2013) 126001.
  doi:10.1088/0029-5515/53/12/126001
\bibitem{rummel2005} Th. Rummel, K. Rieße, H. Ehmler: Manufacture and test of the non-planar coils for
  WENDELSTEIN 7-X. \emph{Fusion Eng.\ Des.} 75--79 (2005) 117--121. doi:10.1016/j.fusengdes.2005.06.118
\bibitem{maix2006} R.\,K. Maix, V. Bagnato, M. Fricke et al.: Design, production and QA test results of the
  NbTi CIC conductors for the W7-X magnet system. \emph{J.\ Phys.: Conf.\ Ser.} 43 (2006) 753--758.
  doi:10.1088/1742-6596/43/1/184
\bibitem{bachmann2024} C. Bachmann, M. Siccinio, A. Ciula et al.: Re-design of EU DEMO with a low aspect ratio.
  \emph{Fusion Eng.\ Des.} 204 (2024) 114518.
\bibitem{gil2025} P.\,F. Gil, V. Brack, T. Schuler, P. Huslage, E.\,V. Stenson: Manufacturing tolerances of
  non-planar coils for an optimized tabletop stellarator. arXiv:2507.22516 (2025).
\bibitem{schmeling2025} D. Schmeling, M. Russo, B.\,T. Gebreamlak et al.: Experimental results from early
  non-planar NI-HTS magnet prototypes for the Columbia Stellarator eXperiment (CSX). \emph{IEEE Trans.\ Plasma
  Sci.} (2026), doi:10.1109/TPS.2026.3681798; arXiv:2510.07042.
\bibitem{simsopt} M. Landreman, B. Medasani, F. Wechsung, A. Giuliani, R. Jorge, C. Zhu: SIMSOPT: A flexible
  framework for stellarator optimization. \emph{J.\ Open Source Softw.} 6(65) (2021) 3525.
  doi:10.21105/joss.03525. W7-X coil data: \texttt{simsopt/configs/W7-X.dat}, coil set \texttt{w7x\_v001}.
\bibitem{desc} D.\,W. Dudt, E. Kolemen: DESC: A stellarator equilibrium solver. \emph{Phys.\ Plasmas} 27 (2020)
  102513.
\bibitem{gmsh} C. Geuzaine, J.-F. Remacle: Gmsh: A 3-D finite element mesh generator with built-in pre- and
  post-processing facilities. \emph{Int.\ J.\ Numer.\ Meth.\ Eng.} 79 (2009) 1309--1331.
\bibitem{elmer} ElmerFEM, version 9.0. CSC -- IT Center for Science, Finland. \url{https://www.elmerfem.org}
\bibitem{freecad} FreeCAD, version 1.1. \url{https://www.freecad.org}
\bibitem{pyvista} C.\,B. Sullivan, A. Kaszynski: PyVista: 3D plotting and mesh analysis through a streamlined
  interface for the Visualization Toolkit (VTK). \emph{J.\ Open Source Softw.} 4(37) (2019) 1450.
\bibitem{paraview} J. Ahrens, B. Geveci, C. Law: ParaView: An end-user tool for large data visualization. In:
  \emph{Visualization Handbook}, Elsevier (2005).
\bibitem{blender} Blender, version 4.0. Blender Foundation. \url{https://www.blender.org}
\bibitem{hermes} Nous Research: Hermes Agent. \url{https://github.com/NousResearch/hermes-agent} (MIT licence).
\bibitem{claudecode} Anthropic: Claude Code. \url{https://www.anthropic.com/claude-code}
\end{thebibliography}
\end{document}
